% p.73 — العوامل ومتعلقاتها (تتمة)
% Layout: twocol. rightcol: entries 132–141, sub-heading اسماء الاشارات, 142–150;
% leftcol: 151–153, the column-width list "الكلمات التي تقع بعدها جملة" (given as two \par
% paragraphs), 154–157.
% Note: reledpar sets the two columns line by line (line i of the right column beside line i of
% the left one, each pair as tall as its taller member), so a column made of tall sarf "lines"
% padded the other column and the page overflowed. Each column's whole content is therefore
% wrapped in one rightbox/leftbox (minipages inside the reledpar columns): the columns then
% flow independently and the page fits; the apparatus numbers each column as a single line.
\mspage{73}
\begin{twocol}
\begin{rightbox}
\begin{sarf}{|c|c|}
(132) اَلَّتِيْ & هٰذِهِ اِمْرَأَةٌ اَلَّتِيْ سَرَقَتْ قَلَمَكَ\\
(133) اَلَّتَانِ & جَاءَتْنِيْ اِمْرَأَتَانِ اَلَّتَانِ جَاءَتَا اَمْسِ\\
 & رَاَيْتُ اِمْرَأَتَيْنِ اَلَّتَيْنِ ضَرَبْتَهُمَا\\
(134) اَلَّاتِيْ & هُنَّ نِسَاءٌ اَلَّاتِيْ اَرْضَعْنَ زَيْدًا\\
(135) اَلَّوَاتِيْ & جَاءَتْ نِسَاءٌ اَلَّوَاتِيْ عَلَّمْنَكَ\\
(136) مَنْ & مَنْ يَضْرِبْنِيْ اَضْرِبْهُ\\
(137) مَا & مَا يَضْرِبْنِيْ اَضْرِبْهُ\\
(138) اَيٌّ & اَيٌّ يُكْرِمْنِيْ اُكْرِمْهُ\\
(139) اَيَّةٌ & اَيَّةُ تَضْرِبْنِيْ اَضْرِبْهَا\\
(140) آلْ & اَلضَّارِبُ ظَالِمٌ\\
 & الضَّارِبَانِ ظَالِمَانِ\\
(141) ذُوْ & \unclear{بِيْرِيْ ذُوْ حَضَرْتُ}\\ % TODO: smudged; the classical example is بِئْرِيْ ذُوْ حَفَرْتُ
\end{sarf}
{\subbab{اسماء الاشارات}}
\begin{sarf}{|c|c|}
(142) هٰذَا & هٰذَا رَجُلٌ\\
(143) هٰذَانِ & هٰذَانِ رَجُلَانِ\\
(144) هٰؤُلَاءِ & هٰؤُلَاءِ رِجَالٌ\\
(145) هٰذِهِ & هٰذِهِ اِمْرَأَةٌ\\
(146) هَاتَانِ & هَاتَانِ اِمْرَأَتَانِ\\
(147) هٰؤُلَاءِ & هٰؤُلَاءِ نِسَاءٌ\\
(148) ذَالِكَ & ذَالِكَ رَجُلٌ\\
(149) ذَانِكَ & ذَانِكَ رَجُلَانِ\\
(150) اُولٰئِكَ & اُولٰئِكَ رِجَالٌ\\
\end{sarf}
\end{rightbox}
\begin{leftbox}
\begin{sarf}{|c|c|}
(151) تِلْكَ & تِلْكَ اِمْرَأَةٌ\\
(152) تَانِكَ & تَانِكَ اِمْرَأَتَانِ\\
(153) اُولٰئِكَ & اُولٰئِكَ نِسَاءٌ\\
\end{sarf}
اَلْكَلِمَاتُ الَّتِيْ تَقَعُ بَعْدَهَا جُمْلَةٌ وَهِيَ
\par
اِنَّمَا اَنَّمَا كَاَنَّمَا لٰكِنَّمَا لَيْتَمَا لَعَلَّمَا بَيْنَمَا بَيْنَا اِذَا اِذْ اِذْمَا اِذَامَا لَوْ لَوْلَا لَوْمَا هَلْ بَلْ كَيْفَ اَيْنَ مَتٰى اِنْ اَنْ لٰكِنْ نَحْوُ اِنَّمَا زَيْدٌ مُنْطَلِقٌ
\par
\begin{sarf}{|c|c|}
(154) اَلظُّرُوْفُ & قَبْلُ بَعْدُ فَوْقُ تَحْتُ\\
 & حَيْثُ اِذَا اِذْ اَيْنَ اَنّٰى\\
 & مَتٰى كَيْفَ اَيَّانَ مُذْ مُنْذُ\\
 & لَدَا يَوْمُ مِثْلُ بَيْنَ مَعَ\\
 & قَطُّ عِوَضُ\\
(155) اَلتَّنَازُعُ & ضَرَبَنِيْ وَاَكْرَمَنِيْ زَيْدٌ\\
 & ضَرَبْتُ وَاَكْرَمْتُ زَيْدًا\\
 & ضَرَبَنِيْ وَاَكْرَمْتُ زَيْدٌ\\
 & ضَرَبْتُ وَاَكْرَمَنِيْ زَيْدًا\\
(156) اَلتَّعَجُّبُ & مَا اَحْسَنَ زَيْدًا\\
 & اَحْسِنْ بِزَيْدٍ\\
(157) اَلْاِشْتِغَالُ & زَيْدًا ضَرَبْتُهُ\\
 & زَيْدٌ ضَرَبْتُ غُلَامَهُ\\
\end{sarf}
\end{leftbox}
\end{twocol}
